% !TEX program = xelatex
\documentclass[12pt,a4paper]{ctexart}
\usepackage[top=2.5cm,bottom=2.5cm,left=2.5cm,right=2.5cm,headheight=14pt]{geometry}
\usepackage{amsmath,amssymb}
\usepackage{enumitem}
\usepackage{booktabs}
\usepackage{longtable}
\usepackage{fancyhdr}
\usepackage{titlesec}
\usepackage{tikz}
\usepackage[table]{xcolor}
\usepackage{tcolorbox}
\tcbuselibrary{skins,breakable}
\usepackage{hyperref}
\usepackage{caption}
\usepackage{fancyvrb}
\usepackage{listings}

\lstset{
  basicstyle=\small\ttfamily,
  breaklines=true,
  frame=single,
  backgroundcolor=\color{gray!5},
  rulecolor=\color{gray!50},
  xleftmargin=2em,
  framexleftmargin=1.5em,
  numbers=left,
  numberstyle=\tiny\color{gray},
}

\titleformat{\section}{\Large\bfseries\centering}{}{0em}{}
\titlespacing{\section}{0pt}{1.5em}{0.8em}
\titleformat{\subsection}{\large\bfseries}{}{0em}{}
\titlespacing{\subsection}{0pt}{1em}{0.5em}
\titleformat{\subsubsection}{\normalsize\bfseries}{}{0em}{}
\titlespacing{\subsubsection}{0pt}{0.8em}{0.3em}

\pagestyle{fancy}
\fancyhf{}
\fancyhead[L]{\small Linux Administration \& Application -- Lab Guide}
\fancyhead[R]{\small Dawei Yuan}
\fancyfoot[C]{\thepage}
\renewcommand{\headrulewidth}{0.4pt}

\newtcolorbox[auto counter]{infobox}{
  colback=blue!3!white,colframe=blue!60!black,fonttitle=\bfseries,title=关键信息~\thetcbcounter,breakable,arc=3mm}
\newtcolorbox{warnbox}{
  colback=orange!5!white,colframe=orange!80!black,fonttitle=\bfseries,title=注意事项,breakable,arc=3mm}
\newtcolorbox{cmdbox}{
  colback=gray!3!white,colframe=gray!60!black,fonttitle=\bfseries,title=命令参考,breakable,arc=3mm}

\renewcommand{\contentsname}{目录}

\begin{document}

% ===================== 封面 =====================
\begin{titlepage}
\begin{center}
\vspace*{3cm}

{\Huge\bfseries Linux Administration \& Application\\[8pt]Comprehensive Lab Guide}\\[14pt]
{\Large Linux管理与应用 --- 综合实训指导书}

\vspace{2cm}

\begin{tabular}{rl}
\textbf{Author:} & Dawei Yuan \\[8pt]
\textbf{Date:} & April 2026 \\[8pt]
\textbf{Platform:} & Rocky Linux 9.8 + VirtualBox
\end{tabular}

\vfill
{\small Version: v2.0 (Rocky Linux 9.8)}
\end{center}
\end{titlepage}

\newpage
\tableofcontents
\newpage

% ===================== 项目列表 =====================
\section{项目（实训）名称}

\begin{enumerate}[label=\arabic*.,leftmargin=*]
  \item YUM配置国内源
  \item 搭建VPN服务器
  \item 搭建邮件服务器
  \item 搭建DHCP和DNS服务器
  \item 搭建Web服务器
  \item 搭建Samba服务器
  \item 搭建代理服务器
  \item 搭建NAT服务器
\end{enumerate}

\section{项目（实训）学时数}

本实训项目预计实训学时数为48课时。

\section{项目（实训）目标}

通过实训，使学生能够完成企业Linux服务器的配置、管理与维护。通过实际操作，使学生掌握一定的操作技能，能认真、细致、准确的操作。通过实践过程，培养学生独立思考、独立工作的能力及团队协作精神。

% ===================== 实训环境准备 =====================
\section{实训环境准备}

\subsection{推荐平台}

\begin{infobox}
  本指导书基于 \textbf{Rocky Linux 9.8} 编写，与 CentOS Stream 9 完全兼容。推荐使用 VirtualBox 7.x 或 VMware Workstation 17 作为虚拟机平台。
\end{infobox}

\subsubsection{为什么选 Rocky Linux 9.8}

\begin{itemize}
  \item Rocky Linux 是 RHEL 的下游发行版，100\% 兼容 RHEL
  \item 9.8 是 9.x 系列最新稳定版（2026-05-25 发布），软件包较新
  \item 社区活跃，被广泛用于服务器和开发环境
  \item 作为 CentOS 的替代品，生态成熟
\end{itemize}

\subsubsection{下载 ISO 镜像}

推荐使用中科大镜像源下载：

\begin{cmdbox}
\small
\begin{lstlisting}[language=bash,numbers=none]
https://mirrors.ustc.edu.cn/rocky/9.8/isos/x86_64/Rocky-9.8-x86_64-boot.iso
\end{lstlisting}
\end{cmdbox}

文件大小约 1.38 GB（Boot ISO），安装时需联网下载所需软件包。

\subsubsection{三种 ISO 对比}

\begin{table}[h]
\centering
\begin{tabular}{p{2cm}p{2.5cm}p{3cm}p{5cm}}
\toprule
ISO 类型 & 文件大小 & 安装方式 & 适用场景 \\
\midrule
Boot ISO  & 约1.4 GB & 网络安装，按需下载 & \textbf{推荐}—实验、虚拟机测试 \\
DVD ISO   & 约14 GB  & 完整离线安装    & 无网络环境、批量部署 \\
Minimal ISO & 约2.6 GB & 离线最小化安装  & 纯命令行服务器 \\
\bottomrule
\end{tabular}
\end{table}

\subsection{VirtualBox 安装}

\subsubsection{下载 VirtualBox}

\begin{table}[h]
\centering
\begin{tabular}{p{3cm}p{10cm}}
\toprule
来源       & 下载地址 \\
\midrule
官方       & \url{https://www.virtualbox.org/wiki/Downloads} \\
中科大镜像  & \url{https://mirrors.ustc.edu.cn/virtualbox/} \\
\bottomrule
\end{tabular}
\end{table}

选择对应系统的安装包下载。Windows 系统下载文件名类似 \texttt{VirtualBox-7.0.24-167081-Win.exe}。

\subsubsection{安装 VirtualBox}

\begin{enumerate}[label=\arabic*.,leftmargin=*]
  \item 双击安装包，进入安装向导
  \item 按默认选项完成安装（网卡驱动、USB 驱动均勾选）
  \item 安装完成后，桌面出现 Oracle VM VirtualBox 图标
\end{enumerate}

\subsection{VirtualBox 虚拟机配置}

\begin{table}[h]
\centering
\begin{tabular}{ll}
\toprule
配置项  & 建议值 \\
\midrule
类型    & Linux / Red Hat (64-bit) \\
内存    & 2048 MB 以上 \\
CPU     & 2 核以上 \\
虚拟硬盘 & 20 GB 动态分配 \\
网络    & \textbf{NAT + Host-Only 双网卡}（外网+固定IP SSH） \\
显存    & 128 MB（如需 GUI） \\
\bottomrule
\end{tabular}
\end{table}

\subsubsection{创建步骤}

\begin{enumerate}[label=\arabic*.,leftmargin=*]
  \item VirtualBox → 新建
  \item 名称：\texttt{Rocky-9.8}，类型：\texttt{Linux}，版本：\texttt{Red Hat (64-bit)}
  \item 内存：建议 2048 MB
  \item 虚拟硬盘：现在创建虚拟硬盘 → VDI → 动态分配 → 20 GB
  \item 创建完成后 → 设置 → 存储 → 选择 Boot ISO 镜像
  \item 网络设置：\textbf{网卡1：NAT} + \textbf{网卡2：Host-Only}（混杂模式均选"拒绝"）
\end{enumerate}

\subsubsection{安装建议}

\begin{itemize}
  \item 安装两台虚拟机：一台作为服务器（2GB 内存+2核CPU），一台作为客户端
  \item 软件选择建议 Minimal Install（纯命令行），体量最小、组件干净
  \item 安装时在 Installation Destination 选 Automatic 自动分区即可
\end{itemize}

\subsubsection{克隆第二台虚拟机（客户端）}

实验2（VPN）、实验4（DHCP+DNS）、实验6（Samba）、实验8（NAT）需要两台虚拟机协同实验。使用 VirtualBox 克隆功能可快速生成第二台虚拟机，无需重新安装系统。

\begin{enumerate}[label=\arabic*.,leftmargin=*]
  \item \textbf{克隆操作}
  \begin{enumerate}[label=(\arabic*)]
    \item 关闭原虚拟机（关机，不是挂起）
    \item VirtualBox 主界面 → 右键原虚拟机 → \textbf{Clone}
    \item 名称：\texttt{Rocky-9.8-client}，路径默认
    \item \textbf{勾选「重新初始化所有网卡的 MAC 地址」}（防止 IP 冲突）
    \item 克隆类型选 \textbf{完全复制}（独立不依赖原盘）
  \end{enumerate}

  \item \textbf{克隆后配置}
\end{enumerate}

\begin{cmdbox}
\small
\begin{lstlisting}[language=bash,numbers=none]
# 1. 修改主机名
sudo hostnamectl set-hostname client

# 2. 修复 NAT 网卡（MAC 地址变化后需重新创建连接）
sudo nmcli connection add type ethernet ifname enp0s8 con-name nat autoconnect yes

# 3. 重新生成 SSH 主机密钥（避免与原机冲突）
sudo rm -f /etc/ssh/ssh_host_*
sudo ssh-keygen -A
sudo systemctl restart sshd
\end{lstlisting}
\end{cmdbox}

\begin{enumerate}[label=\arabic*.,leftmargin=*]
  \setcounter{enumi}{2}
  \item \textbf{验证网络互通}
\end{enumerate}

克隆完成后，两台虚拟机的 Host-Only 网卡在同一网段，可互通：

\begin{center}
\begin{tabular}{lccc}
\toprule
虚拟机 & 主机名 & Host-Only IP & 角色 \\
\midrule
原机器 & server & 192.168.56.101 & 承载服务 \\
克隆机 & client & 192.168.56.102 & 测试客户端 \\
\bottomrule
\end{tabular}
\end{center}

\begin{cmdbox}
\small
\begin{lstlisting}[language=bash,numbers=none]
# 在两台 VM 上分别查看 IP
ip a | grep 192.168.56

# 互相 ping 测试
ping 192.168.56.101    # 在 client 上执行
ping 192.168.56.102    # 在 server 上执行
\end{lstlisting}
\end{cmdbox}

\subsubsection{安装后初始配置}

\begin{cmdbox}
\small
\begin{lstlisting}[language=bash,numbers=none]
sudo hostnamectl set-hostname server
yum update -y
yum install -y vim git curl wget net-tools bind-utils
\end{lstlisting}
\end{cmdbox}

\subsection{SSH 远程访问配置}

为方便从 Windows 终端远程操作 Linux 虚拟机，推荐配置 SSH 服务和 NAT + Host-Only 双网卡方案。

\subsubsection{启用 SSH 服务}

\begin{cmdbox}
\small
\begin{lstlisting}[language=bash,numbers=none]
yum install -y openssh-server
systemctl enable sshd --now
firewall-cmd --permanent --add-service=ssh
firewall-cmd --reload
\end{lstlisting}
\end{cmdbox}

\subsubsection{网络模式选型}

\begin{table}[h]
\centering
\begin{tabular}{p{2.2cm}p{4cm}p{3.5cm}p{3.5cm}}
\toprule
模式       & 原理 & Windows→虚拟机SSH & 虚拟机上网 \\
\midrule
\textbf{NAT + Host-Only} & \textbf{双网卡分工} & \textbf{✓ 固定IP直接连} & \textbf{✓} \\
桥接       & 接入物理局域网  & ✓ 直接连（WiFi不稳定） & ✓ \\
NAT 单卡   & 藏在宿主机NAT后 & 需端口转发 & ✓ \\
Host-Only 单卡 & 虚拟独立局域网  & ✓ 直接连 & ✗ \\
\bottomrule
\end{tabular}
\end{table}

\begin{infobox}
  \textbf{推荐方案：NAT + Host-Only 双网卡} \\
  网卡1 (NAT)：虚拟机上外网（yum install、curl） \\
  网卡2 (Host-Only)：固定 IP 192.168.56.x，Windows SSH 连接，不受 WiFi 切换影响 \\
  Linux 内核自动按目的 IP 选择走哪张网卡，无需手动配置路由。 \\
  \textbf{WiFi 下不推荐桥接模式}：WiFi 桥接不支持完整 ARP/DHCP 透传，常导致只获得 IPv6 地址、IPv4 不可用。
\end{infobox}

\subsubsection{双网卡配置步骤}

\begin{enumerate}[label=\arabic*.,leftmargin=*]
  \item 确保虚拟机已关机（不是挂起），进入 VirtualBox → 设置 → 网络
  \item 网卡1：启用网络连接 ☑ → 连接方式：\textbf{NAT} → 混杂模式：拒绝
  \item 网卡2：启用网络连接 ☑ → 连接方式：\textbf{仅主机(Host-Only)网络} → 混杂模式：拒绝
  \item 启动虚拟机，查看 Host-Only 网卡 IP：
  \begin{cmdbox}
  \begin{lstlisting}[language=bash,numbers=none]
ip a | grep 192.168.56
  \end{lstlisting}
  \end{cmdbox}
  \item 在 Windows 终端（PowerShell / CMD）连接：
  \begin{cmdbox}
  \begin{lstlisting}[language=bash,numbers=none]
ssh 用户名@192.168.56.101
  \end{lstlisting}
  \end{cmdbox}
\end{enumerate}

\begin{warnbox}
  每个网卡的\textbf{启用网络连接}必须勾选，\textbf{混杂模式}全部选"拒绝"。\\
  确定按钮变灰通常是虚拟机未关机，需先关闭电源再进设置。
\end{warnbox}

\subsection{Guest Additions 安装（可选）}

Guest Additions 提供剪贴板共享、拖拽文件、自适应分辨率等功能。

\begin{cmdbox}
\small
\begin{lstlisting}[language=bash,numbers=none]
yum install -y gcc kernel-devel kernel-headers elfutils-libelf-devel
# VirtualBox 菜单：设备 → 安装增强功能
mkdir -p /mnt/cdrom
mount /dev/cdrom /mnt/cdrom
cd /mnt/cdrom && ./VBoxLinuxAdditions.run
reboot
\end{lstlisting}
\end{cmdbox}

重启后通过 VirtualBox 菜单 → 设备 → 共享剪贴板 → 双向，即可开启复制粘贴。

% ===================== 实验1 =====================
\newpage
\section{子项目1：YUM配置国内源}

\subsection{实验目的}

将 Rocky Linux 9.8 默认的国外 YUM 源替换为阿里云镜像源，提升软件包下载速度。

\subsection{背景说明}

Rocky Linux 的软件包管理工具 YUM 默认配置的是国外官方源，在国内网络环境下往往会导致软件包下载速度缓慢甚至失败。通过将 YUM 源更换为国内主流镜像源，可显著提升软件包下载速度。

\subsection{操作步骤}

\begin{enumerate}[label=\arabic*.,leftmargin=*]

  \item \textbf{备份原有配置}

  \begin{cmdbox}
  \begin{lstlisting}[language=bash]
mkdir -p /etc/yum.repos.d/backup
cp /etc/yum.repos.d/*.repo /etc/yum.repos.d/backup/
  \end{lstlisting}
  \end{cmdbox}

  \item \textbf{配置阿里云镜像源}

  在 \texttt{/etc/yum.repos.d/} 目录下创建 \texttt{rocky-aliyun.repo} 文件：

  \begin{cmdbox}
  \begin{lstlisting}[language=bash]
cat > /etc/yum.repos.d/rocky-aliyun.repo << 'EOF'
[baseos-aliyun]
name=Rocky Linux $releasever - BaseOS (Aliyun)
baseurl=https://mirrors.aliyun.com/rockylinux/$releasever/BaseOS/$basearch/os/
gpgcheck=1
enabled=1
gpgkey=file:///etc/pki/rpm-gpg/RPM-GPG-KEY-Rocky-9

[appstream-aliyun]
name=Rocky Linux $releasever - AppStream (Aliyun)
baseurl=https://mirrors.aliyun.com/rockylinux/$releasever/AppStream/$basearch/os/
gpgcheck=1
enabled=1
gpgkey=file:///etc/pki/rpm-gpg/RPM-GPG-KEY-Rocky-9

[extras-aliyun]
name=Rocky Linux $releasever - Extras (Aliyun)
baseurl=https://mirrors.aliyun.com/rockylinux/$releasever/extras/$basearch/os/
gpgcheck=1
enabled=1
gpgkey=file:///etc/pki/rpm-gpg/RPM-GPG-KEY-Rocky-9
EOF
  \end{lstlisting}
  \end{cmdbox}

  \item \textbf{清理缓存并重建}

  \begin{cmdbox}
  \begin{lstlisting}[language=bash]
yum clean all
yum makecache
  \end{lstlisting}
  \end{cmdbox}

  \item \textbf{验证}

  \begin{cmdbox}
  \begin{lstlisting}[language=bash]
yum repolist        # 查看启用仓库列表
yum search vim      # 测试搜索功能
  \end{lstlisting}
  \end{cmdbox}

\end{enumerate}

\subsection{验收脚本}

\begin{cmdbox}
\small
\begin{lstlisting}[language=bash]
sudo bash labs/lab1_yum_repo/setup.sh   # 一键配置
sudo bash labs/lab1_yum_repo/verify.sh  # 验收检查
\end{lstlisting}
\end{cmdbox}

\subsection{验收标准}

\begin{itemize}
  \item rocky-aliyun.repo 配置文件存在
  \item yum repolist 显示阿里云仓库已启用
  \item yum makecache 正常完成
  \item yum search 可正常使用
\end{itemize}

% ===================== 实验2 =====================
\newpage
\section{子项目2：搭建VPN服务器}

\subsection{实验目的}

使用 OpenVPN 搭建 VPN 服务器，实现客户端通过加密隧道安全访问内网资源。

\subsection{背景说明}

虚拟私人网络（Virtual Private Network，简称 VPN）是一种用于连接中、大型企业或团体间私人网络的通讯方法。VPN 的讯息透过公用的网络架构（例如互联网）来传送内网的网络讯息。本次实训使用 OpenVPN 完成 VPN 服务器的搭建。

\textbf{原理简述}：OpenVPN 在服务端和客户端之间建立一条加密的虚拟隧道（TUN 模式，三层 IP 隧道）。双方各创建一个 tun0 虚拟网卡，分配 VPN 子网地址（10.8.0.0/24）。实际数据包经由物理网卡发出前，被 OpenVPN 加密并封装为 UDP 报文；到达对端后解密还原，注入对端 tun0。服务端通过 \texttt{push "route"} 将内网路由推送给客户端，使客户端能访问服务端所在的内网资源。

\subsection{操作步骤}

\begin{enumerate}[label=\arabic*.,leftmargin=*]

  \item \textbf{安装 OpenVPN 和 easy-rsa}

  \begin{cmdbox}
  \begin{lstlisting}[language=bash]
yum install -y epel-release
yum install -y openvpn easy-rsa wget
  \end{lstlisting}
  \end{cmdbox}

  \item \textbf{配置 easy-rsa 变量}

  \begin{cmdbox}
  \begin{lstlisting}[language=bash]
mkdir -p /opt/easy-rsa && cd /opt/easy-rsa
cp -a /usr/share/easy-rsa/3/* ./

cat > vars << 'EOF'
if [ -z "$EASYRSA_CALLER" ]; then
    echo "You appear to be sourcing an Easy-RSA 'vars' file." >&2
    return 1
fi
set_var EASYRSA_DN "cn_only"
set_var EASYRSA_REQ_COUNTRY "CN"
set_var EASYRSA_REQ_PROVINCE "Beijing"
set_var EASYRSA_REQ_CITY "Beijing"
set_var EASYRSA_REQ_ORG "example"
set_var EASYRSA_REQ_EMAIL "admin@example.com"
set_var EASYRSA_NS_SUPPORT "yes"
EOF
  \end{lstlisting}
  \end{cmdbox}

  \item \textbf{生成证书}

  \begin{cmdbox}
  \begin{lstlisting}[language=bash]
./easyrsa init-pki
echo "" | ./easyrsa build-ca nopass
echo "" | ./easyrsa gen-req server nopass
./easyrsa --batch sign-req server server
openssl ecparam -genkey -name prime256v1 -out /opt/easy-rsa/pki/ec.pem
echo "" | ./easyrsa gen-req client nopass
./easyrsa --batch sign-req client client
  \end{lstlisting}
  \end{cmdbox}

  \item \textbf{服务端配置}（\texttt{/etc/openvpn/server/server.conf}）

  \begin{cmdbox}
  \begin{lstlisting}[language=bash]
mkdir -p /etc/openvpn/server
cat > /etc/openvpn/server/server.conf << 'EOF'
port 1194
proto udp
dev tun
ca /opt/easy-rsa/pki/ca.crt
cert /opt/easy-rsa/pki/issued/server.crt
key /opt/easy-rsa/pki/private/server.key
dh none
ecdh-curve prime256v1
server 10.8.0.0 255.255.255.0
push "route 192.168.56.0 255.255.255.0"
keepalive 10 120
max-clients 100
status openvpn-status.log
log /var/log/openvpn.log
verb 3
client-to-client
persist-key
persist-tun
duplicate-cn
comp-lzo
EOF
  \end{lstlisting}
  \end{cmdbox}

  \item \textbf{启动服务端}（含旧配置清理和证书开放）

  \begin{cmdbox}
  \begin{lstlisting}[language=bash]
# 清理旧配置（支持重复执行）
systemctl stop openvpn-server@server 2>/dev/null || true
systemctl disable openvpn-server@server 2>/dev/null || true
rm -rf /opt/easy-rsa/pki /etc/openvpn/server

firewall-cmd --permanent --add-port=1194/udp
firewall-cmd --reload
systemctl enable openvpn-server@server
systemctl start openvpn-server@server

# 开放证书读权限，使客户端可用同名用户 scp 获取
chmod -R o+rX /opt/easy-rsa/pki
  \end{lstlisting}
  \end{cmdbox}

\end{enumerate}

\subsection{客户端配置}

客户端需要从服务端获取3个证书文件，然后编写 client.conf 并启动连接。

\begin{enumerate}[label=\arabic*.,leftmargin=*]

  \item \textbf{获取证书}

  \begin{cmdbox}
  \begin{lstlisting}[language=bash]
# 自动检测当前用户名（兼容 sudo 场景）
SSH_USER="${SUDO_USER:-$USER}"
SERVER=192.168.56.101

yum install -y openvpn
mkdir -p /etc/openvpn/client
scp ${SSH_USER}@${SERVER}:/opt/easy-rsa/pki/ca.crt           /etc/openvpn/client/
scp ${SSH_USER}@${SERVER}:/opt/easy-rsa/pki/issued/client.crt /etc/openvpn/client/
scp ${SSH_USER}@${SERVER}:/opt/easy-rsa/pki/private/client.key /etc/openvpn/client/
  \end{lstlisting}
  \end{cmdbox}

  \item \textbf{客户端配置文件}（\texttt{/etc/openvpn/client/client.conf}）

  \begin{cmdbox}
  \begin{lstlisting}[language=bash]
cat > /etc/openvpn/client/client.conf << EOF
client
dev tun
proto udp
remote <server_ip> 1194
ca /etc/openvpn/client/ca.crt
cert /etc/openvpn/client/client.crt
key /etc/openvpn/client/client.key
resolv-retry infinite
nobind
persist-key
persist-tun
verb 3
comp-lzo
pull-filter ignore "route 192.168.56.0"
EOF
  \end{lstlisting}
  \end{cmdbox}

  \item \textbf{启动客户端}

  \begin{cmdbox}
  \begin{lstlisting}[language=bash]
systemctl enable openvpn-client@client
systemctl start openvpn-client@client
  \end{lstlisting}
  \end{cmdbox}

\end{enumerate}

\subsection{验证}

\subsubsection{服务端验收}

\begin{cmdbox}
\begin{lstlisting}[language=bash]
sudo bash labs/lab2_VPN/server_verify.sh
\end{lstlisting}
\end{cmdbox}

\begin{itemize}
  \item openvpn-server@server 正在运行
  \item 1194/udp 端口在监听
  \item tun0 网卡已创建
  \item server.conf 配置文件存在
\end{itemize}

\subsubsection{客户端验收}

\begin{cmdbox}
\begin{lstlisting}[language=bash]
sudo bash labs/lab2_VPN/client_verify.sh
\end{lstlisting}
\end{cmdbox}

\begin{itemize}
  \item openvpn-client@client 正在运行
  \item tun0 已创建并分配 VPN IP（10.8.0.x）
  \item 可 ping 通服务端 VPN 地址 10.8.0.1
  \item 路由表存在 tun0 相关路由，推送已生效
\end{itemize}

\subsection{停止 VPN 服务}

实验完成后，在两台虚拟机上分别停止 OpenVPN 服务，释放 tun0 虚拟网卡：

\subsubsection{服务端}

\begin{cmdbox}
\begin{lstlisting}[language=bash]
sudo systemctl stop openvpn-server@server
sudo systemctl disable openvpn-server@server
\end{lstlisting}
\end{cmdbox}

\subsubsection{客户端}

\begin{cmdbox}
\begin{lstlisting}[language=bash]
sudo systemctl stop openvpn-client@client
sudo systemctl disable openvpn-client@client
\end{lstlisting}
\end{cmdbox}

停止后 \texttt{ip a} 中 tun0 消失，Host-Only 网络恢复正常。

% ===================== 实验3 =====================
\newpage
\section{子项目3：搭建邮件服务器}

\subsection{实验目的}

使用 Postfix 搭建 SMTP 邮件发送客户端，通过 QQ 邮箱 SMTP 中继实现邮件发送功能。

\subsection{背景说明}

\textbf{本实验并非搭建完整公网邮件服务器。} 完整邮件系统需 Postfix + Dovecot + 自有域名 + MX DNS 记录。本实验中的 Postfix 仅作为 SMTP 中继客户端，通过 QQ 邮箱 SMTP 服务器代理发送邮件。

\textbf{工作流程}：Postfix 接收本地邮件 → SASL 认证登录 \texttt{smtp.qq.com:587}（需 QQ 邮箱 + 授权码）→ QQ 验证身份后以 QQ 名义将邮件投递到最终收件人。

Postfix 作为工业级 MTA，通过纯文本配置文件即可操控完整 SMTP 协议行为：\texttt{relayhost} 指定中继目标、\texttt{smtp\_use\_tls} 控制传输加密、\texttt{smtp\_sasl\_auth\_enable} 启用身份认证。\textbf{配置文件 → 程序解析 → 协议行为变化}这一范式贯穿全部 8 个实验，是本课程的核心教学模式。

\subsection{操作步骤}

\begin{enumerate}[label=\arabic*.,leftmargin=*]

  \item \textbf{安装 Postfix}

  \begin{cmdbox}
  \begin{lstlisting}[language=bash]
yum install -y postfix mailx cyrus-sasl-plain
  \end{lstlisting}
  \end{cmdbox}

  \item \textbf{配置 Postfix main.cf}

  在 \texttt{/etc/postfix/main.cf} 末尾添加：

  \begin{cmdbox}
  \begin{lstlisting}[language=bash]
relayhost = [smtp.qq.com]:587
smtp_sasl_auth_enable = yes
smtp_sasl_password_maps = hash:/etc/postfix/sasl_passwd
smtp_sasl_security_options = noanonymous
smtp_use_tls = yes
smtp_tls_security_level = encrypt
smtp_tls_CAfile = /etc/pki/tls/certs/ca-bundle.crt
  \end{lstlisting}
  \end{cmdbox}

  \item \textbf{创建认证文件}

  创建 \texttt{/etc/postfix/sasl\_passwd}：

  \begin{cmdbox}
  \begin{lstlisting}[language=bash]
echo "[smtp.qq.com]:587 your_qq@qq.com:授权码" > /etc/postfix/sasl_passwd
chmod 600 /etc/postfix/sasl_passwd
postmap /etc/postfix/sasl_passwd
  \end{lstlisting}
  \end{cmdbox}

  \begin{warnbox}
    授权码需要在 QQ 邮箱→设置→账户→POP3/SMTP服务 中生成，不是 QQ 密码。
  \end{warnbox}

  \item \textbf{启动并测试}

  \begin{cmdbox}
  \begin{lstlisting}[language=bash]
systemctl enable postfix --now
echo "Test mail from Rocky Linux" | mail -s "Test Subject" recipient@qq.com
  \end{lstlisting}
  \end{cmdbox}

  \item \textbf{问题排查}

  如果收不到邮件，查看日志：

  \begin{cmdbox}
  \begin{lstlisting}[language=bash]
tail -f /var/log/maillog
  \end{lstlisting}
  \end{cmdbox}

\end{enumerate}

\subsection{验收脚本}

\begin{cmdbox}
\begin{lstlisting}[language=bash]
sudo bash labs/lab3_email/setup.sh    # 会提示输入 QQ 邮箱和授权码
sudo bash labs/lab3_email/verify.sh
\end{lstlisting}
\end{cmdbox}

\subsection{验收标准}

\begin{itemize}
  \item postfix 服务正在运行
  \item main.cf 中 relayhost 已配置为 QQ SMTP
  \item /etc/postfix/sasl\_passwd 和 .db 文件存在
  \item postfix check 语法检查通过
  \item 发送测试邮件成功
\end{itemize}

% ===================== 实验4 =====================
\newpage
\section{子项目4：搭建DHCP和DNS服务器}

\subsection{实验目的}

使用 dnsmasq 一站式搭建 DHCP 和 DNS 服务器，实现局域网 IP 地址动态分配和域名解析。

\subsection{操作步骤}

\begin{enumerate}[label=\arabic*.,leftmargin=*]

  \item \textbf{安装 dnsmasq}

  \begin{cmdbox}
  \begin{lstlisting}[language=bash]
yum install -y dnsmasq
  \end{lstlisting}
  \end{cmdbox}

  \item \textbf{配置 dnsmasq}

  编辑 \texttt{/etc/dnsmasq.conf}：

  \begin{cmdbox}
  \small
  \begin{lstlisting}[language=bash]
# DNS 配置
listen-address=127.0.0.1,192.168.56.10
bind-interfaces
domain=tecmint.lan
server=223.5.5.5        # 阿里 DNS
server=223.6.6.6        # 阿里 DNS
address=/tecmint.lan/192.168.56.10

# DHCP 配置
dhcp-range=192.168.56.50,192.168.56.150,12h
dhcp-option=3,192.168.56.1
dhcp-option=6,192.168.56.10
  \end{lstlisting}
  \end{cmdbox}

  \item \textbf{配置防火墙}

  \begin{cmdbox}
  \begin{lstlisting}[language=bash]
firewall-cmd --permanent --add-service=dns
firewall-cmd --permanent --add-service=dhcp
firewall-cmd --reload
  \end{lstlisting}
  \end{cmdbox}

  \item \textbf{启动服务}

  \begin{cmdbox}
  \begin{lstlisting}[language=bash]
systemctl enable dnsmasq --now
dnsmasq --test    # 测试配置是否正确
  \end{lstlisting}
  \end{cmdbox}

  \item \textbf{验证 DNS}

  \begin{cmdbox}
  \begin{lstlisting}[language=bash]
nslookup tecmint.lan 127.0.0.1
nslookup www.baidu.com 127.0.0.1
  \end{lstlisting}
  \end{cmdbox}

  \item \textbf{验证 DHCP}

  在客户端虚拟机上将网卡设置为 DHCP 模式，确保与 DHCP 服务器在同一个 Host-Only 网络中，连接后检查是否获取到 IP 地址。

\end{enumerate}

\subsection{验收脚本}

\begin{cmdbox}
\begin{lstlisting}[language=bash]
sudo bash labs/lab4_DHCP_DNS/setup.sh
sudo bash labs/lab4_DHCP_DNS/verify.sh
\end{lstlisting}
\end{cmdbox}

\subsection{验收标准}

\begin{itemize}
  \item dnsmasq 服务正在运行
  \item DHCP 地址池已配置
  \item DNS 域名 tecmint.lan 已配置
  \item nslookup 可解析 tecmint.lan
  \item 配置语法正确
\end{itemize}

% ===================== 实验5 =====================
\newpage
\section{子项目5：搭建Web服务器}

\subsection{实验目的}

使用 Apache (httpd) 搭建 Web 服务器，实现基于域名的虚拟主机，使一个 IP 地址承载多个网站。

\subsection{操作步骤}

\begin{enumerate}[label=\arabic*.,leftmargin=*]

  \item \textbf{安装 Apache 和 BIND}

  \begin{cmdbox}
  \begin{lstlisting}[language=bash]
yum install -y httpd bind
  \end{lstlisting}
  \end{cmdbox}

  \item \textbf{创建虚拟主机目录和默认页面}

  \begin{cmdbox}
  \begin{lstlisting}[language=bash]
mkdir -p /var/www/test-web1 /var/www/test-web2
echo "this is web1" > /var/www/test-web1/index.html
echo "this is web2" > /var/www/test-web2/index.html
  \end{lstlisting}
  \end{cmdbox}

  \item \textbf{配置 DNS 区域解析}

  在 \texttt{/etc/named.conf} 末尾添加两个 zone 定义：

  \begin{cmdbox}
  \small
  \begin{lstlisting}[language=bash]
zone "test-web1.com" IN { type master; file "web1.com.zone"; };
zone "test-web2.com" IN { type master; file "web2.com.zone"; };
  \end{lstlisting}
  \end{cmdbox}

  创建区域文件 \texttt{/var/named/web1.com.zone} 和 \texttt{/var/named/web2.com.zone}，内容中将
  \texttt{www} 的 A 记录指向服务器 IP。

  \item \textbf{配置虚拟主机}

  创建 \texttt{/etc/httpd/conf.d/vhosts.conf}：

  \begin{cmdbox}
  \small
  \begin{lstlisting}[language=bash]
<VirtualHost *:80>
    DocumentRoot /var/www/test-web1
    ServerName www.test-web1.com
    DirectoryIndex index.html
</VirtualHost>

<VirtualHost *:80>
    DocumentRoot /var/www/test-web2
    ServerName www.test-web2.com
    DirectoryIndex index.html
</VirtualHost>
  \end{lstlisting}
  \end{cmdbox}

  \item \textbf{启动服务并验证}

  \begin{cmdbox}
  \begin{lstlisting}[language=bash]
systemctl enable named httpd
systemctl restart named httpd

# 测试虚拟主机
curl -H "Host: www.test-web1.com" http://127.0.0.1/
curl -H "Host: www.test-web2.com" http://127.0.0.1/
  \end{lstlisting}
  \end{cmdbox}

\end{enumerate}

\subsection{验收脚本}

\begin{cmdbox}
\begin{lstlisting}[language=bash]
sudo bash labs/lab5_Web/setup.sh
sudo bash labs/lab5_Web/verify.sh
\end{lstlisting}
\end{cmdbox}

\subsection{验收标准}

\begin{itemize}
  \item httpd 和 named 服务正在运行
  \item curl test-web1.com 返回``this is web1''
  \item curl test-web2.com 返回``this is web2''
  \item 虚拟主机配置文件存在
\end{itemize}

% ===================== 实验6 =====================
\newpage
\section{子项目6：搭建Samba服务器}

\subsection{实验目的}

使用 Samba 搭建跨平台文件共享服务器，实现多用户、多组、分权限的文件共享方案。

\subsection{业务场景}

该公司分支机构有 system、develop、productdesign 和 test 共 4 个小组。个人办公机操作系统为 Windows/Linux，服务器操作系统为 Rocky Linux 9.8。需要设计一套安全文件共享方案。

\subsection{权限矩阵}

\bgroup\small
\begin{table}[h]
\centering
\begin{tabular}{lcccccc}
\toprule
目录        & 所有者 & develop & productdesign & test & develop\_test & 匿名 \\
\midrule
/data/share      & system    & ---      & ---           & ---  & ---   & --- \\
/data/share/develop & develop  & 读写     & 不可访问       & 不可访问 & --- & --- \\
/data/share/productdesign & productdesign & 不可访问 & 读写   & 不可访问 & ---   & --- \\
/data/share/test & test      & 不可访问  & 不可访问        & 读写   & ---   & --- \\
/data/share/library & system  & 只读     & 只读          & 只读   & ---   & 只读 \\
/data/share/develop\_test & --- & 读写    & 不可访问        & 读写   & 读写   & --- \\
/data/share/temp & system    & 读写     & 读写          & 读写   & ---   & 读写 \\
\bottomrule
\end{tabular}
\caption{Samba 共享权限矩阵}
\end{table}
\egroup

\subsection{操作步骤}

\begin{enumerate}[label=\arabic*.,leftmargin=*]

  \item \textbf{安装 Samba}

  \begin{cmdbox}
  \begin{lstlisting}[language=bash]
yum install -y samba samba-client samba-common cifs-utils
  \end{lstlisting}
  \end{cmdbox}

  \item \textbf{创建目录结构}

  \begin{cmdbox}
  \begin{lstlisting}[language=bash]
mkdir -p /data/share/{develop,productdesign,test,library,develop_test,temp}
  \end{lstlisting}
  \end{cmdbox}

  \item \textbf{创建用户和组}

  \begin{cmdbox}
  \begin{lstlisting}[language=bash]
groupadd -f system
groupadd -f develop
groupadd -f productdesign
groupadd -f test
groupadd -f develop_test

useradd -g develop  -G develop_test -d /data/share/develop  -s /sbin/nologin develop
useradd -g productdesign -G develop_test -d /data/share/productdesign -s /sbin/nologin productdesign
useradd -g test -G develop_test -d /data/share/test -s /sbin/nologin test
useradd -g system -G develop,productdesign,test,develop_test -d /data/share -s /sbin/nologin system
  \end{lstlisting}
  \end{cmdbox}

  \item \textbf{设置目录权限}

  \begin{cmdbox}
  \begin{lstlisting}[language=bash]
chmod 755  /data/share
chown system:system /data/share
chmod 2770 /data/share/develop /data/share/productdesign /data/share/test /data/share/develop_test
chmod 3777 /data/share/temp
chmod 755  /data/share/library

chown develop:system      /data/share/develop
chown productdesign:system /data/share/productdesign
chown test:system         /data/share/test
chown system:system       /data/share/library
chown system:system       /data/share/temp

setfacl -m g:develop:rwx /data/share/develop_test
setfacl -m g:test:rwx    /data/share/develop_test
  \end{lstlisting}
  \end{cmdbox}

  \item \textbf{配置 /etc/samba/smb.conf}

  \begin{cmdbox}
  \small
  \begin{lstlisting}[language=bash]
[global]
    workgroup = SYSTEM
    security = user
    map to guest = Bad User

[develop]
    path = /data/share/develop
    writeable = yes
    valid users = develop,@system

[productdesign]
    path = /data/share/productdesign
    writeable = yes
    valid users = productdesign,@system

[test]
    path = /data/share/test
    writeable = yes
    valid users = test,@system

[library]
    path = /data/share/library
    writeable = no
    guest ok = yes

[develop_test]
    path = /data/share/develop_test
    writeable = yes
    valid users = system,@develop_test

[temp]
    path = /data/share/temp
    writeable = yes
    guest ok = yes
  \end{lstlisting}
  \end{cmdbox}

  \item \textbf{添加 Samba 用户}

  \begin{cmdbox}
  \begin{lstlisting}[language=bash]
smbpasswd -a system     # 输入密码
smbpasswd -a develop    # 输入密码
smbpasswd -a productdesign  # 输入密码
smbpasswd -a test       # 输入密码

testparm -v     # 测试配置
systemctl enable smb --now
  \end{lstlisting}
  \end{cmdbox}

\end{enumerate}

\subsection{验收脚本}

\begin{cmdbox}
\begin{lstlisting}[language=bash]
sudo bash labs/lab6_Samba/setup.sh
sudo bash labs/lab6_Samba/verify.sh
\end{lstlisting}
\end{cmdbox}

\begin{warnbox}
  setup.sh 中所有 Samba 用户默认密码为 \texttt{123456}。正式环境中请修改。
\end{warnbox}

\subsection{验收标准}

\begin{itemize}
  \item smb 服务正在运行
  \item testparm 配置语法正确
  \item 7 个共享均已定义
  \item 各用户目录存在且权限正确
  \item develop 用户可访问 develop 共享，被拒绝访问 productdesign 共享
\end{itemize}

% ===================== 实验7 =====================
\newpage
\section{子项目7：搭建代理服务器}

\subsection{实验目的}

使用 Squid 搭建 HTTP 正向代理服务器，实现客户端通过代理访问外部网络。

\subsection{操作步骤}

\begin{enumerate}[label=\arabic*.,leftmargin=*]

  \item \textbf{安装 Squid}

  \begin{cmdbox}
  \begin{lstlisting}[language=bash]
yum install -y squid
  \end{lstlisting}
  \end{cmdbox}

  \item \textbf{配置代理}

  编辑 \texttt{/etc/squid/squid.conf}，设置代理端口 8080 并允许内网访问：

  \begin{cmdbox}
  \begin{lstlisting}[language=bash]
http_port 8080
acl localnet src 192.168.0.0/16
http_access allow localnet
  \end{lstlisting}
  \end{cmdbox}

  \item \textbf{启动并测试}

  \begin{cmdbox}
  \begin{lstlisting}[language=bash]
systemctl enable squid --now

# 通过代理访问网页
curl -x http://localhost:8080 http://www.baidu.com
  \end{lstlisting}
  \end{cmdbox}

  \item \textbf{客户端使用}

  在客户端的浏览器或终端中配置代理地址 \texttt{http://<服务器IP>:8080} 即可。

\end{enumerate}

\subsection{验收脚本}

\begin{cmdbox}
\begin{lstlisting}[language=bash]
sudo bash labs/lab7_proxy/setup.sh
sudo bash labs/lab7_proxy/verify.sh
\end{lstlisting}
\end{cmdbox}

\subsection{验收标准}

\begin{itemize}
  \item squid 服务正在运行
  \item 8080 端口在监听
  \item curl -x 通过代理可访问外网（HTTP 200/301/302）
\end{itemize}

% ===================== 实验8 =====================
\newpage
\section{子项目8：搭建NAT服务器}

\subsection{实验目的}

使用 iptables 配置 NAT（网络地址转换）服务器，使内网主机通过该服务器访问外网，同时配置防火墙规则限制访问。

\subsection{背景说明}

NAT（Network Address Translation，网络地址转换）用于在本地网络使用私有地址、连接互联网时转换为全局 IP 地址的技术。本次实训使用 iptables 作为 NAT 服务器接入网络。

\subsection{前置条件}

服务器需要至少两张网卡：一张外网（eth0/NAT）、一张内网（eth1/Host-Only）。

\subsection{操作步骤}

\begin{enumerate}[label=\arabic*.,leftmargin=*]

  \item \textbf{安装 iptables-services}

  \begin{cmdbox}
  \begin{lstlisting}[language=bash]
yum install -y iptables-services
  \end{lstlisting}
  \end{cmdbox}

  \item \textbf{启用 IP 转发}

  \begin{cmdbox}
  \begin{lstlisting}[language=bash]
echo "net.ipv4.ip_forward = 1" >> /etc/sysctl.conf
sysctl -w net.ipv4.ip_forward=1
  \end{lstlisting}
  \end{cmdbox}

  \item \textbf{停止 firewalld，启用 iptables}

  \begin{cmdbox}
  \begin{lstlisting}[language=bash]
systemctl stop firewalld
systemctl disable firewalld
systemctl enable iptables --now
  \end{lstlisting}
  \end{cmdbox}

  \item \textbf{配置 NAT 规则}

  \begin{cmdbox}
  \begin{lstlisting}[language=bash]
# 清空现有规则
iptables -F; iptables -t nat -F; iptables -X; iptables -t nat -X

# 默认策略
iptables -P INPUT DROP
iptables -P FORWARD DROP
iptables -P OUTPUT ACCEPT

# 允许回环和已建立连接
iptables -A INPUT  -i lo -j ACCEPT
iptables -A INPUT  -m state --state ESTABLISHED,RELATED -j ACCEPT
iptables -A FORWARD -m state --state ESTABLISHED,RELATED -j ACCEPT

# 允许 SSH
iptables -A INPUT -p tcp --dport 22 -j ACCEPT

# NAT 转发：内网→外网
iptables -A FORWARD -i eth1 -o eth0 -j ACCEPT

# 只允许 Web、DNS、邮件服务
iptables -A FORWARD -i eth0 -o eth1 -p tcp --dport 80 -j ACCEPT
iptables -A FORWARD -i eth0 -o eth1 -p tcp --dport 443 -j ACCEPT
iptables -A FORWARD -i eth0 -o eth1 -p udp --dport 53 -j ACCEPT
iptables -A FORWARD -i eth0 -o eth1 -p tcp --dport 25 -j ACCEPT
iptables -A FORWARD -i eth0 -o eth1 -p tcp --dport 110 -j ACCEPT

# SNAT 地址伪装
iptables -t nat -A POSTROUTING -o eth0 -j MASQUERADE

# 保存规则
iptables-save > /etc/sysconfig/iptables
  \end{lstlisting}
  \end{cmdbox}

  \item \textbf{验证规则}

  \begin{cmdbox}
  \begin{lstlisting}[language=bash]
iptables -L -v -n           # 查看 filter 表
iptables -t nat -L -v -n    # 查看 nat 表
  \end{lstlisting}
  \end{cmdbox}

\end{enumerate}

\subsection{验收脚本}

\begin{cmdbox}
\begin{lstlisting}[language=bash]
sudo bash labs/lab8_NAT/setup.sh
sudo bash labs/lab8_NAT/verify.sh
\end{lstlisting}
\end{cmdbox}

\subsection{验收标准}

\begin{itemize}
  \item net.ipv4.ip\_forward = 1
  \item iptables 服务正在运行
  \item FORWARD 链存在 ACCEPT 规则
  \item POSTROUTING 链存在 MASQUERADE 规则
  \item 规则已持久化到 /etc/sysconfig/iptables
  \item firewalld 已停止（避免冲突）
\end{itemize}

% ===================== 验收脚本说明 =====================
\newpage
\section{验收脚本使用说明}

每个实验目录下均提供了两个脚本：

\begin{table}[h]
\centering
\begin{tabular}{lp{10cm}}
\toprule
文件 & 作用 \\
\midrule
\texttt{setup.sh}  & 一键安装配置脚本，学生执行以完成实验 \\
\texttt{verify.sh} & 验收检查脚本，教师用于检查学生实验是否正确完成 \\
\bottomrule
\end{tabular}
\end{table}

所有 \texttt{verify.sh} 输出统一格式：

\begin{cmdbox}
\begin{lstlisting}[language=bash,numbers=none]
[PASS] 检查项描述
[FAIL] 检查项描述
[WARN] 检查项描述
=== 结果：N 通过, M 失败 ===
\end{lstlisting}
\end{cmdbox}

\begin{itemize}
  \item 全部 PASS → exit 0（验收通过）
  \item 有 FAIL → exit 1（验收不通过）
\end{itemize}

批量检查所有实验：

\begin{cmdbox}
\begin{lstlisting}[language=bash]
for lab in labs/lab*; do
    echo "=== Checking $(basename $lab) ==="
    sudo bash "$lab/verify.sh" || echo "FAILED"
    echo ""
done
\end{lstlisting}
\end{cmdbox}

% ===================== 考核标准 =====================
\section{考核或评价标准}

\begin{table}[h]
\centering
\begin{tabular}{p{3cm}p{8cm}p{2cm}}
\toprule
考核项目     & 考核方法                                   & 分值 \\
\midrule
上机操作情况 & 根据课堂考勤、课堂表现、调试表现综合评定     & 40   \\
实训报告质量 & 撰写准确、美观，能完整清晰分析实训结果        & 30   \\
答辩情况     & 答辩时能正确回答问题                         & 30   \\
\midrule
合计         &                                            & 100  \\
\bottomrule
\end{tabular}
\end{table}

% ===================== 进度安排 =====================
\section{进度安排}

\begin{table}[h]
\centering
\begin{tabular}{clc}
\toprule
序号 & 实训内容              & 课时安排 \\
\midrule
1  & YUM配置国内源          & 2课时  \\
2  & 搭建邮件服务器          & 6课时  \\
3  & 搭建VPN服务器           & 6课时  \\
4  & 搭建DHCP和DNS服务器     & 6课时  \\
5  & 搭建Web服务器           & 6课时  \\
6  & 搭建Samba服务器         & 6课时  \\
7  & 搭建代理服务器          & 6课时  \\
8  & 搭建NAT服务器           & 6课时  \\
9  & 答辩                    & 4课时  \\
\midrule
   & 合计                    & 48课时 \\
\bottomrule
\end{tabular}
\end{table}

% ===================== 辅助资料 =====================
\section{辅助教学资料}

\begin{itemize}
  \item \url{http://www.linux.org} — The Linux Home Page at Linux Online
  \item \url{http://bbs.chinaunix.net} — 中国最大的Linux/Unix技术社区
  \item \url{https://docs.rockylinux.org} — Rocky Linux 官方文档
  \item \url{https://github.com/davidyuan666/gdust-linux-labs} — 本课程实训代码仓库（含全部 setup.sh 和 verify.sh）
\end{itemize}

\end{document}
